% answers_short.tex —— 2006 年·上海春季招生数学试卷 · 简短答案速查
% 由 tools/make_short_answers.py 自 answers.tex 的 \gksolhead/\gkanswer 提取，
% 并与 2006/2006_数学真题_春季上海_答案要点.md 逐题比对后生成，勿手改（改源 MD 后重跑）。

\begin{gknotice}
\textbf{使用说明}：本节只给最终答案便于快速核对，完整推导见解析版。答案为编写者独立推导，\textbf{非官方标准答案}。
\end{gknotice}

\gkpart{填空题}
\begin{itemize}[leftmargin=2.6em]
  \item[\textbf{1.}] $\dfrac34$
  \item[\textbf{2.}] $2$
  \item[\textbf{3.}] $\dfrac{x-5}{3}$，$x\in[5,8]$
  \item[\textbf{4.}] $\left(-1,\dfrac12\right)$
  \item[\textbf{5.}] $(0,\sqrt{10})$
  \item[\textbf{6.}] $-x-x^4$
  \item[\textbf{7.}] $48$ 种
  \item[\textbf{8.}] $\dfrac{16}{3}$
  \item[\textbf{9.}] $\dfrac7{25}$
  \item[\textbf{10.}] $2$
  \item[\textbf{11.}] $4$
  \item[\textbf{12.}] $\frac{a_1+\cdots+a_m}{m}\leq\frac{a_1+\cdots+a_n}{n}\leq\frac{a_{m+1}+\cdots+a_n}{n-m}\quad(1\leq m<n)$
\end{itemize}

\gkpart{选择题}
\begin{center}\begin{tabular}{ccccc}
\toprule
题号 & 13 & 14 & 15 & 16 \\
\midrule
答案 & B & C & A & B \\
\bottomrule
\end{tabular}\end{center}

\gkpart{解答题}
\begin{itemize}[leftmargin=2.6em]
  \item[\textbf{17.}] $\arccos\dfrac9{25}$
  \item[\textbf{18.}] $x^2-6x+10=0$
  \item[\textbf{19.}] (1) $\dfrac{4\sqrt3+3}{5}$；(2) $[1,2]$
  \item[\textbf{20.}] (1) $y=-\dfrac17x^2+\dfrac{64}{7}$；(2) $|AC|=2\sqrt5$，$|BC|=4$
  \item[\textbf{21.}] (1) 按分段表达式绘图；(2) $B\subsetneq A=(-\infty,2-\sqrt{14}]\cup[0,4]\cup[2+\sqrt{14},+\infty)$；(3) 见解析版
  \item[\textbf{22.}] (1) $d=3$；(2) $a_{30}=10(1+d+d^2)$，值域 $[\dfrac{15}{2},+\infty)$；(3) $a_{10(n+1)}=10(1+d+\cdots+d^n)$。固定 $n$ 为奇数时值域 $\mathbb R\setminus\{10\}$；固定 $n$ 为偶数时，令 $x_n>0$ 为 $(n+1)x_n^n+n x_n^{n+1}=1$ 的唯一正根，$\mu_n=\dfrac{1+x_n^{n+1}}{1+x_n}$，值域为 $[10\mu_n,+\infty)$。$n=2$ 时 $x_2=\dfrac12$、$\mu_2=\dfrac34$，与（2）相同
\end{itemize}
